\documentclass[12pt]{article}
\usepackage{authblk}

% Math packages
\usepackage{amsmath,amsthm,amssymb,amsfonts}
\usepackage{physics} % Clean vector, derivative, norm, etc.

% Graphics and tables
\usepackage{graphicx}
\usepackage{float} % Improved float handling
\usepackage{booktabs} % Nicer tables
\usepackage{lscape}
\usepackage{pdflscape}
\usepackage{framed}
\usepackage{pifont}
\usepackage{array}
\usepackage{enumitem}
\usepackage[most]{tcolorbox}

\providecommand{\RR}{\mathbb{R}}

\usepackage{tikz}
\usetikzlibrary{arrows.meta, positioning}

% Hyperlinks
\usepackage{hyperref}

% Optional: define styling once
\tcbset{
  highlight theorem/.style={
    colback=blue!5!white,
    colframe=blue!75!black,
    fonttitle=\bfseries,
    coltitle=black,
    boxrule=0.8pt,
    arc=3pt,
    title={\textbf{Classical Implication Template: $P$ vs.\ $NP$ (Conditional)}},
    before upper={\parindent15pt},
  }
}

% Define blank tape symbol
\newcommand{\blank}{\sqcup}


% ============================
% Complexity class and WCC macro definitions
% ============================
\newcommand{\Pclass}{\mathbf{P}}
\newcommand{\NPclass}{\mathbf{NP}}
\newcommand{\PWCC}{\mathbf{PWCC}}
\newcommand{\NPWCC}{\mathbf{NPWCC}}

% ============================
% Machine model macros
% ============================
\newcommand{\TM}{\mathrm{TM}}
\newcommand{\WCC}{\mathrm{WCC}}

% ============================
% Encoding / decoding macros
% ============================
\newcommand{\enc}{\mathrm{enc}}
\newcommand{\dec}{\mathrm{dec}}

% ============================
% NC hierarchy (small caps)
% ============================
\newcommand{\NC}{\mathsf{NC}}

% ============================
% Misc formatting
% ============================
\newcommand{\SAT}{\mathrm{SAT}}



\title{Discrete Wave-Constrained Computation and Classical Complexity:\\
Turing Equivalence for $P$ and $NP$}
\author{Richard J. Reyes}
\date{Revised October 27, 2026 \\ \small Original Release: November 26, 2025}

\newtheorem{theorem}{Theorem}[section]
\newtheorem{lemma}[theorem]{Lemma}
\newtheorem{corollary}[theorem]{Corollary}
\newtheorem{definition}[theorem]{Definition}
\newtheorem{proposition}[theorem]{Proposition}
\newtheorem{conjecture}[theorem]{Conjecture}
\newtheorem{remark}[theorem]{Remark}
\begin{document}
\maketitle

\begin{abstract}
Wave Confinement Theory (WCT) motivates a discrete model of computation in which
information is stored and processed by local interactions between finitely many
metastable states arranged on a lattice. In this paper we isolate the resulting
\emph{wave-constrained computation} (WCC) model as a finite-state, finite-range
local-update machine and compare its complexity classes to the classical ones.

First, we define a \emph{uniform} discrete WCC machine model with fixed finite
alphabet, bounded-radius local updates, polynomial-size constructible lattices,
and polynomial-time input encoders and output decoders. We explicitly simulate
polynomial-time deterministic Turing machines with constant-factor overhead in
the number of global WCC rounds and polynomial overhead in total work; conversely,
each such WCC machine has a polynomial-time Turing simulation. Thus
$P_{\mathrm{WCC}}^{\mathrm{disc}} = P_{\mathrm{Turing}}$.

We then lift this equivalence from decision procedures to verifier-based
nondeterministic computation. We define $NP_{\mathrm{WCC}}$ via polynomial-time
WCC verifiers with polynomially bounded witnesses, and prove that
$NP_{\mathrm{WCC}} = NP_{\mathrm{Turing}}$.
Taken together, these results show that the uniform discrete abstraction
realizes exactly the classical $P$ and $NP$ classes. Exact symbolic representation
of wave trajectories needs a compatibility condition on quantization and local
evolution. Additional physical curvature, energetic and PDE constraints are
not covered by the discrete complexity-class equalities; the physical
capacity estimates are conditional on their specific dynamical hypotheses.

\end{abstract}














%%%%%%%%%%%%%%%%%%%%%%%%%%%%%%%%%%%%%%%%%%%%%%%%%%%%%%%%%%%%
% 0.x — Resource Model Summary
%%%%%%%%%%%%%%%%%%%%%%%%%%%%%%%%%%%%%%%%%%%%%%%%%%%%%%%%%%%%

\subsection{Resource Model Summary}

For convenience, we summarize the resource model used throughout the
definition of discrete WCC machines. Let $x \in \Sigma^n$ be an
input of length $n$.

\begin{enumerate}

    \item \textbf{Space (Finite Lattice).}
    Computation takes place on a finite lattice
    $\Lambda_n \subset \mathbb{Z}^d$ (with fixed $d$) satisfying
    \[
        |\Lambda_n| \le p(n)
    \]
    for some polynomial $p$. The cardinality $|\Lambda_n|$ represents
    the space cost of the machine.

    \item \textbf{Precision (Finite State).}
    The state set $S$ is finite. Each cell holds an element of $S$,
    representing at most $\ell = \lceil\log_2|S|\rceil$ bits of information.
    No operation increases $\ell$; precision remains fixed and finite
    throughout computation.

    \item \textbf{Locality (Neighborhood Radius).}
    A fixed interaction radius $r \in \mathbb{N}$ defines neighborhoods
    \[
        N_r(i) = \{ j \in \Lambda_n : \|i-j\|_\infty \le r \}.
    \]
    Each update applies a local rule $U : S^{N_r} \to S$.

    \item \textbf{Determinism and Time.}
    Computation proceeds in synchronous, deterministic rounds via $U$.
    For a WCC machine $W$, the time cost on inputs of length $n$ is
    \[
        T_W(n) =
        \max_{x \in \Sigma^n}
        \min\{ t\ge0 : \mathrm{dec}(C_t) \in \{0,1\}\}.
    \]
    Given $C_0$ and $U$, the entire trajectory $(C_t)_{t \ge 0}$ is uniquely
    determined (deterministic replay).

    \item \textbf{Total Configuration Space.}
    Since $|\Lambda_n|$ and $|S|$ are finite,
    the global configuration space $S^{\Lambda_n}$ is finite.
    This precludes hypercomputation via infinite-dimensional state spaces.

\end{enumerate}

The machine descriptions, lattice enumerators, encoders, and decoders are
fixed for a language and execute in polynomial time. This uniformity is
essential: otherwise an unrestricted input encoding could decide the language
before the first WCC update.

%%%%%%%%%%%%%%%%%%%%%%%%%%%%%%%%%%%%%%%%%%%%%%%%%%%%%%%%%%%%
% Lemma — Precision Stability
%%%%%%%%%%%%%%%%%%%%%%%%%%%%%%%%%%%%%%%%%%%%%%%%%%%%%%%%%%%%

\begin{lemma}[Precision Stability]
\label{lem:precision-stability}
Let
$W = (\{\Lambda_n\}, S, r, U, \mathrm{enc}, \mathrm{dec})$
be a discrete WCC machine.
If $C_t : \Lambda_n \to S$ is representable using
$\ell = \lceil\log_2 |S|\rceil$ bits per cell, then so is $C_{t+1}$ obtained by the
local update rule~$U$.
\end{lemma}

\begin{proof}
For each $i \in \Lambda_n$,
\[
   C_{t+1}(i) = U(C_t|_{N_r(i)}).
\]
Since $U : S^{N_r} \to S$ is a total function whose image lies in the
same finite state set $S$, its outputs require no more than
$\lceil\log_2|S|\rceil$ bits.
No arithmetic involving unbounded precision is performed.
Thus, the precision requirement remains constant, preventing
accumulation of infinite information during evolution.
\end{proof}

%%%%%%%%%%%%%%%%%%%%%%%%%%%%%%%%%%%%%%%%%%%%%%%%%%%%%%%%%%%%
% Proposition — Local Implementability of TM Transitions
%%%%%%%%%%%%%%%%%%%%%%%%%%%%%%%%%%%%%%%%%%%%%%%%%%%%%%%%%%%%

\begin{proposition}[Local Implementability of TM Transitions]
\label{prop:local-tm}
Let $M$ be a deterministic single-tape Turing machine with transition
function
\[
   \delta : Q \times \Gamma \to Q \times \Gamma \times \{L, R\}.
\]
Under the two-phase head/symbol encoding of Section~1, each transition
of $M$ is realizable by a WCC local update rule with radius $r = 1$.
\end{proposition}

\begin{proof}
In the simulation, each tape cell is encoded as a WCC state
$(\gamma, q^\#, p) \in S$, where $\gamma \in \Gamma$,
$q^\# \in Q \cup \{\bot, (q,L), (q,R)\}$ encodes the head state or a
directional mark, and $p \in \{0,1\}$ is a phase bit.

To realize a TM transition, the rule $U$ need only inspect the cell
currently holding the head and its immediate neighbors (to which the
head may move in the next phase).  These cells lie inside the radius--1
neighborhood $\{i-1, i, i+1\}$.

Given $(q, \gamma)$ at the head location, $U$ rewrites the symbol,
enters the next state, and places a directional mark for the subsequent
phase.  Movement of the head by one cell is therefore computed using
only local information. Hence, all TM transitions are locally
implementable in the WCC model.
\end{proof}

%%%%%%%%%%%%%%%%%%%%%%%%%%%%%%%%%%%%%%%%%%%%%%%%%%%%%%%%%%%%
% Proposition — Deterministic Evolution and Replay
%%%%%%%%%%%%%%%%%%%%%%%%%%%%%%%%%%%%%%%%%%%%%%%%%%%%%%%%%%%%

\begin{proposition}[Deterministic Evolution and Replay]
\label{prop:deterministic-replay}
Let $W$ be a discrete WCC machine.
Then the global evolution
\[
   C_{t+1}(i) = U(C_t|_{N_r(i)})
\]
is deterministic. Consequently, given $C_0$ and $U$, the entire
trajectory $(C_t)_{t \ge 0}$ is uniquely determined and can be
reconstructed (replayed) by any verifier.
\end{proposition}

\begin{proof}
Since $U$ is a function from the finite domain
$S^{N_r}$ to the finite set $S$, each cell update is uniquely
determined by the configuration of its neighborhood. Thus $C_{t+1}$ is
uniquely determined by $C_t$.

By induction, the full sequence of configurations is uniquely
determined by $C_0$. A verifier with access to $C_0$ and $U$ can
therefore reproduce the entire evolution without ambiguity,
establishing deterministic replay.
\end{proof}


% ============================
% Discrete WCC model and Turing equivalence
% ============================

\section{Discrete WCC Machines and Turing Equivalence}
\label{sec:tm_wcc_equiv}

In this section we isolate the discrete core of wave-constrained computation
(WCC) as a finite-state local-update machine model and prove that its
polynomial-time class coincides with the classical deterministic polynomial-time
class. This makes precise the identification
\[
P_{\mathrm{WCC}} = P_{\mathrm{Turing}}
\]
used elsewhere in the paper.

We first define the discrete WCC machine under the locality and precision rails,
then construct an explicit simulation of any polynomial-time Turing machine by a
polynomial-time WCC machine with only constant-factor overhead in time. We then
recall the converse simulation (WCC by Turing machines) and obtain full
equivalence.

\subsection{Discrete WCC machine model}

Let $\Sigma$ be a finite input alphabet and let $n = |x|$ denote input length.

% WIN-09: enforce polynomial uniformity of lattice, encoding and decoding.
\begin{definition}[Uniform discrete WCC machine]
\label{def:wcc_machine}
A \emph{uniform discrete WCC machine} is a tuple
\[
 W=(\{\Lambda_n\}_{n\ge1},S,r,U,\mathrm{enc},\mathrm{dec})
\]
with the following data, fixed independently of an individual input.
\begin{enumerate}
\item The spatial dimension $d\ge1$ and neighborhood radius $r\ge0$
are constants. The cell set $\Lambda_n\subset\mathbb{Z}^d$ satisfies
$|\Lambda_n|\le p(n)$ and lies in a box with polynomially bounded
integer coordinates. A fixed deterministic algorithm constructs and
enumerates $\Lambda_n$ and its fixed-radius neighborhoods in time
$\mathrm{poly}(n)$.
\item $S$ is a fixed, finite alphabet. Boundary cells outside $\Lambda_n$
read a distinguished fixed background symbol $s_\partial\in S$.
\item $U:S^{\{-r,\ldots,r\}^d}\to S$ is a single fixed local rule
of constant description length. Its synchronous action is
\[
 C_{t+1}(i)=U\bigl((\overline C_t(i+v))_{\|v\|_\infty\le r}\bigr),
 \quad i\in\Lambda_n,
\]
where $\overline C_t$ extends $C_t$ by the background symbol.
\item The uniform encoder $\mathrm{enc}$ constructs
$C_0=\mathrm{enc}(x)\in S^{\Lambda_{|x|}}$ in time
$\mathrm{poly}(|x|)$ on a standard deterministic Turing machine.
\item The uniform decoder $\mathrm{dec}_n:S^{\Lambda_n}\to
\{\mathsf{run},0,1\}$ is computed in time $\mathrm{poly}(n)$
from the canonical configuration encoding. It returns $\mathsf{run}$
before halting and $0$ (reject) or $1$ (accept) on a halting
configuration. Halting outputs are persistent (equivalently the
computation is stopped at its first halting configuration).
\end{enumerate}
The polynomial bounds, fixed rule, and algorithms are part of the
machine description, not extra advice indexed by $n$ or an oracle.
All complexity classes denoted $P_{\mathrm{WCC}}$ and
$NP_{\mathrm{WCC}}$ below refer to this \emph{discrete uniform} model
unless a physical superscript is explicitly attached.
\end{definition}



\begin{definition}[Polynomial-time discrete WCC]
\label{def:pwcc}
For a uniform discrete WCC machine $W$, set
\[
 T_W(n)=\max_{x\in\Sigma^n}\min\{t\ge0:
 \mathrm{dec}_n(C_t)\in\{0,1\}\}.
\]
We require the minimum to exist for every input. The machine decides
$L\subseteq\Sigma^*$ in polynomial time if $T_W(n)\le p(n)$ for
some polynomial $p$ and its first halting output equals
$\mathbf{1}_L(x)$ for every $x$. By Definition~\ref{def:wcc_machine},
the total cost of constructing, simulating, and decoding each of these
runs is polynomial in $n$ whenever the round count is polynomial.
We write $P_{\mathrm{WCC}}^{\mathrm{disc}}$, or simply
$P_{\mathrm{WCC}}$ in the discrete sections, for the resulting class.
\end{definition}



The remaining WCC rails (coercivity, entropy pruning, renewability) impose
additional structure on how such a discrete machine can arise from an underlying
wave-evolution model. In this section we restrict attention to the discrete
interface; the physical realization is addressed separately
(Remark~\ref{rem:wct_realization}).
















% WIN-10: exact symbolic factor requires quantization compatibility.
\subsection{Wave--Symbol Correspondence Under Compatibility}
\label{subsec:wave-symbol-isomorphism}

The local symbolic update $U$ can faithfully describe a discretized
wave evolution only if its quantization commutes with that evolution.
Locality and finite precision alone do not imply this compatibility.
Here a local wave update is a map
$F:Z^{\{-r,\dots,r\}^d}\to Z$ on a finite-precision
subset $Z\subset\mathbb C$, acting on $\Psi_t:\Lambda_n\to Z$ by
\begin{equation}
\Psi_{t+1}(i)=F\bigl((\overline\Psi_t(i+v))_{\|v\|_\infty\le r}\bigr).
\label{eq:psi-local-update}
\end{equation}
This is a discrete local map; realizing it by a particular WCT PDE
requires independent physical analysis.

\begin{definition}[Quantization and representative encoding]
Let $S$ be a fixed finite alphabet. Let
$Q:Z\to S$ be surjective and $\widetilde Q:S\to Z$
select a representative for each symbol, so that
\[
Q\circ\widetilde Q=\mathrm{id}_S.
\]
Both maps act cellwise on configurations. The symbolic state associated
with a field is $C_t=Q\circ\Psi_t$.
\end{definition}

Given $F,Q,\widetilde Q$, define the candidate symbolic rule
\begin{equation}
U(c)=Q\bigl(F(\widetilde Q(c))\bigr),\qquad
c\in S^{\{-r,\ldots,r\}^d}.
\label{eq:U-from-F}
\end{equation}

\begin{theorem}[Exact wave--symbol correspondence]
\label{thm:wave-symbol-isomorphism}
Suppose the actual field trajectory obeys \eqref{eq:psi-local-update}
and that, on every neighborhood configuration $z$ reached by this
trajectory, the \emph{quantization compatibility} identity holds:
\begin{equation}
Q(F(z))=Q\bigl(F(\widetilde Q(Q(z)))\bigr).
\label{eq:quantization-compatibility}
\end{equation}
Then $C_t=Q\circ\Psi_t$ satisfies the exact symbolic update
\begin{equation}
C_{t+1}(i)=U\bigl((\overline C_t(i+v))_{\|v\|_\infty\le r}\bigr).
\label{eq:symbolic-update-from-psi}
\end{equation}
Conversely, for every symbolic local rule $U$ there exists an
\emph{abstract} finite-precision representative evolution, defined by
\[
F_U(z)=\widetilde Q\bigl(U(Q(z))\bigr),
\]
for which every symbolic trajectory is realized exactly on the invariant
code manifold $\widetilde Q(S)^{\Lambda_n}$. This converse is a statement
about a constructed finite-state local evolution, not about solutions of
an independently fixed continuous WCT PDE.
\end{theorem}

\begin{proof}
For an actual trajectory and each site $i$, write $z$ for the
corresponding input neighborhood. By compatibility,
\[
Q(\Psi_{t+1}(i))=Q(F(z))
=Q(F(\widetilde Q(Q(z))))=U(Q(z)),
\]
which is \eqref{eq:symbolic-update-from-psi}.
For the converse, initialize $\Psi_0=\widetilde Q(C_0)$ and use $F_U$.
If $\Psi_t=\widetilde Q(C_t)$, then
\[
\Psi_{t+1}=\widetilde Q\bigl(U(Q(\Psi_t))\bigr)
=\widetilde Q(U(C_t))=\widetilde Q(C_{t+1}).
\]
Induction proves exact representative dynamics and invariance of the
code manifold. No continuous curvature or energy bound is inferred.
\end{proof}

\begin{remark}[Compatibility is necessary]
\label{rem:wave-symbol-obstruction}
Without \eqref{eq:quantization-compatibility}, quantization need not
commute with a local update. For example, let
$Z=\{-1,\tfrac12,1\}$ and $S=\{0,1\}$, with
$Q(z)=\mathbf{1}_{\{z>0\}}$, $\widetilde Q(0)=-1$,
$\widetilde Q(1)=1$. Define the one-cell map
$F(-1)=-1$, $F(\tfrac12)=\tfrac12$, and $F(1)=-1$.
For $\Psi_0=\tfrac12$, one has $Q(F(\Psi_0))=1$,
whereas $U(Q(\Psi_0))=Q(F(1))=0$.
Thus the unqualified wave--symbol isomorphism is false, even for
single-site local maps. Exact correspondence requires the stated
compatibility or an invariant representative code manifold.
\end{remark}


\subsection{Gandy Correspondence and Physical Cellular Automata}
\label{subsec:gandy-correspondence}

Gandy's physical Church--Turing thesis characterizes a broad class of
``reasonable'' physical computation models by four structural axioms:
bounded information density, spatial homogeneity, bounded propagation
speed, and local causality. In this subsection we show that discrete WCC
machines satisfy these axioms and therefore instantiate a physical cellular
automaton in Gandy's sense.

\begin{definition}[Gandy-style physical axioms]
A discrete dynamical system on configurations $C : \Lambda \to S$ satisfies
Gandy's physical constraints if:
\begin{enumerate}
  \item \textbf{Bounded information density:} $S$ is finite, so each site
    carries at most $\log_2 |S|$ bits of information.
  \item \textbf{Spatial homogeneity:} There exists a fixed local rule
    $U : S^{N_r} \to S$ independent of the site $i \in \Lambda$ and of
    the configuration size.
  \item \textbf{Bounded propagation:} There is a fixed radius $r \in \mathbb{N}$
    such that the update at site $i$ depends only on the values in the
    neighbourhood $N_r(i)$.
  \item \textbf{Local causality:} The global update $C \mapsto C'$ is defined
    pointwise by
    \[
    C'(i) = U\big(C|_{N_r(i)}\big),
    \qquad i \in \Lambda,
    \]
    with no dependence on nonlocal aggregates.
\end{enumerate}
\end{definition}

\begin{proposition}[WCC satisfies Gandy axioms]
\label{prop:gandy-wcc}
Let
\[
W = (\{\Lambda_n\}_{n \ge 1}, S, r, U, \mathrm{enc}, \mathrm{dec})
\]
be a discrete WCC machine as in Definition~\ref{def:wcc_machine}.
Then for each $n$ the induced evolution on configurations
$C_t : \Lambda_n \to S$ satisfies all four Gandy-style physical axioms.
\end{proposition}

\begin{proof}
Fix $n \ge 1$.

(1) \emph{Bounded information density.} By definition of a WCC machine, the
state set $S$ is finite. Thus each cell carries at most $\log_2|S|$ bits of
information.

(2) \emph{Spatial homogeneity.} The local update rule $U : S^{N_r} \to S$
is fixed once and for all in the specification of $W$, and does not depend
on $i \in \Lambda_n$ or on $n$.

(3) \emph{Bounded propagation.} The neighbourhood shape $N_r(i)$ is defined
with a fixed radius $r \in \mathbb{N}$. Therefore the update at site $i$
cannot depend on configuration values outside $N_r(i)$, and information
cannot propagate faster than $r$ cells per global step.

(4) \emph{Local causality.} By definition of the WCC evolution,
\[
C_{t+1}(i) = U\big(C_t|_{N_r(i)}\big),
\]
so the future value at site $i$ depends only on the local neighbourhood
$N_r(i)$ at time $t$ and not on any global property of $C_t$.

Thus all four axioms hold for the evolution on $\Lambda_n$, for each $n$.
\end{proof}

\begin{corollary}[Gandy--WCC correspondence]
\label{cor:gandy-wcc-correspondence}
Every discrete WCC machine induces, for each input length $n$, a
finite-state cellular-automaton-style evolution satisfying the listed
structural locality axioms. This is a statement about the abstract
machine, not a proof of realizability by any fixed physical PDE.
\end{corollary}

Proposition~\ref{prop:gandy-wcc} and
Corollary~\ref{cor:gandy-wcc-correspondence} show that uniform WCC is a discrete abstraction compatible with the
structural locality axioms. Physical feasibility and energy accounting are
independent questions.





\subsection{Deterministic Turing machines}

We recall the standard deterministic Turing machine model and its
polynomial-time class.

\begin{definition}[Deterministic Turing machine]
\label{def:tm}
A deterministic Turing machine is a tuple
\[
M = (Q, \Gamma, \Sigma, \delta, q_0, q_{\mathrm{acc}}, q_{\mathrm{rej}})
\]
where:
\begin{itemize}
  \item $Q$ is a finite set of states.
  \item $\Gamma$ is a finite tape alphabet containing a distinguished blank
    symbol $\blank$.
  \item $\Sigma \subseteq \Gamma \setminus \{\blank\}$ is the input alphabet.
  \item $\delta : Q \times \Gamma \to Q \times \Gamma \times \{L,R\}$ is the
    transition function.
  \item $q_0 \in Q$ is the start state.
  \item $q_{\mathrm{acc}}, q_{\mathrm{rej}} \in Q$ are halting accept/reject
    states.
\end{itemize}
A \emph{configuration} of $M$ is a triple $(\gamma, h, q)$ consisting of a
bi-infinite tape $\gamma : \mathbb{Z} \to \Gamma$ with all but finitely many
entries equal to $\blank$, a head position $h \in \mathbb{Z}$, and a current
state $q \in Q$. One transition step applies $\delta(q,\gamma(h))$ and updates
$(\gamma,h,q)$ accordingly.
\end{definition}

\begin{definition}[Polynomial-time Turing computation]
\label{def:pturing}
Let $T_M(n)$ denote the worst-case number of transition steps that $M$ uses on
inputs of length $n$. A language $L \subseteq \Sigma^*$ is in $P_{\mathrm{Turing}}$
if there exists a machine $M$ and constants $c,k \ge 1$ such that
\[
T_M(n) \le c n^k
\]
for all $n$, and $M$ accepts $x$ if and only if $x \in L$.
\end{definition}

\subsection{Simulation of polynomial-time Turing machines by WCC}

We now construct, for each polynomial-time Turing machine $M$, a discrete WCC
machine $W_M$ that simulates $M$ with only constant-factor overhead in time. The
construction proceeds via a one-dimensional radius-$1$ local update rule that
encodes the tape and head in per-cell symbols and uses a two-phase schedule to
implement the Turing transition function.

\subsubsection{Configuration encoding}

Fix a Turing machine $M$ as in Definition~\ref{def:tm}. We first define the
local alphabet used by the WCC machine.

\begin{definition}[Local tape state and phase]
\label{def:local_state}
Let
\[
S_{\mathrm{tape}} := \Gamma \times (Q \cup \{\bot\})
\]
where the symbol $\bot$ denotes ``no head present''. A local tape state
$(\gamma,q) \in S_{\mathrm{tape}}$ thus records the tape symbol $\gamma \in
\Gamma$ and, if $q \in Q$, the presence of the head with state $q$ at the
corresponding cell.

We extend the state space with a phase bit and a directional mark. Let
\[
Q' := Q \cup \{(q,L) : q \in Q\} \cup \{(q,R) : q \in Q\}
\]
and define
\[
S := \Gamma \times (Q' \cup \{\bot\}) \times \{0,1\}.
\]
An element $s = (\gamma, q^\#, p) \in S$ consists of:
\begin{itemize}
  \item a symbol $\gamma \in \Gamma$,
  \item either $q^\# = \bot$ (no head), or $q^\# \in Q$ (head present with
    state $q^\#$), or $q^\# = (q,d)$ with $q \in Q$ and $d \in \{L,R\}$
    indicating that in the next phase the head should move one cell in
    direction $d$,
  \item a phase bit $p \in \{0,1\}$. Phase $p=0$ corresponds to ``before Turing
    transition'' and phase $p=1$ to ``after symbol rewrite, before head
    movement''.
\end{itemize}
\end{definition}

Given a Turing configuration $(\gamma, h, q)$ at time $t$, we encode it as a
configuration $C_t : \mathbb{Z} \to S$ by
\[
C_t(i) :=
\begin{cases}
(\gamma(i), q, 0), & \text{if } i = h,\\[2pt]
(\gamma(i), \bot, 0), & \text{if } i \neq h.
\end{cases}
\]

In the WCC machine we restrict to a finite window
\[
\Lambda_n := \{-L(n), \dots, L(n)\} \subset \mathbb{Z}
\]
with $L(n)$ chosen as a polynomial upper bound on the space usage of $M$ on
inputs of length $n$, for example
\[
L(n) := T_M(n) + n.
\]

\subsubsection{Local update rule}

We now define the local rule $U : S^{\{-1,0,1\}} \to S$, which acts with
radius $r = 1$ in dimension $d=1$. For each site $i \in \Lambda_n$, the
neighbourhood is
\[
N_1(i) = \{i-1,i,i+1\}
\]
(with the understanding that we only care about neighbourhoods fully contained
in $\Lambda_n$; since the head never reaches the boundary within $T_M(n)$
steps, the behaviour of $U$ at the outermost sites can be defined arbitrarily).

The action of $U$ depends on the phase bit of the centre cell.

\paragraph{Phase 0: symbol rewrite and direction tagging.}

In phase $p=0$, at most one cell carries a head state in $Q$, all other cells
have $q^\# = \bot$. For a cell at site $i$ with local state
$s_i = (\gamma_i, q_i^\#, 0)$, we define the phase-$0$ component
$U_{\mathrm{phase0}}$ as:
\begin{itemize}
  \item If $q_i^\# \in Q$ (the head is present at $i$ with state $q_i^\#$), we
    consult the Turing transition function
    \[
    \delta(q_i^\#, \gamma_i) = (q', \gamma', d), \qquad d \in \{L,R\},
    \]
    and set
    \[
    U_{\mathrm{phase0}}(s_i) := (\gamma', (q',d), 1),
    \]
    that is, we overwrite the symbol with $\gamma'$, store a directional mark
    $(q',d)$, and flip the phase bit to $1$.
  \item If $q_i^\# = \bot$ (no head), then the cell does not change its symbol
    or head component and only flips the phase:
    \[
    U_{\mathrm{phase0}}(s_i) := (\gamma_i, \bot, 1).
    \]
\end{itemize}
This part of the rule is radius-$0$; it will be incorporated into the full
radius-$1$ rule below.

\paragraph{Phase 1: head movement.}

In phase $p=1$, exactly one cell in a valid configuration carries a directional
mark $(q',d) \in Q'$ and all others carry either $\bot$ or an unmarked $q \in
Q$. The head must move one cell in direction $d$, and the directional mark must
be cleared.

The update for phase $1$ uses the radius-$1$ neighbourhood
\[
(C(i-1), C(i), C(i+1)) = (s_{i-1}, s_i, s_{i+1}) \in S^3.
\]
We write $s_j = (\gamma_j, q_j^\#, 1)$ with $q_j^\# \in Q' \cup \{\bot\}$ and
define $U_{\mathrm{phase1}} : S^3 \to S$ by case distinction:

\begin{itemize}
  \item If the \emph{centre} cell carries a right-moving mark, i.e.~$q_i^\# =
    (q',R)$, then the head is leaving site $i$:
    \[
    U_{\mathrm{phase1}}(s_{i-1}, s_i, s_{i+1}) := (\gamma_i, \bot, 0).
    \]
  \item If the \emph{left neighbour} carries a right-moving mark and the centre
    cell is headless, i.e.~$q_{i-1}^\# = (q',R)$ and $q_i^\# = \bot$, then the
    head arrives from the left:
    \[
    U_{\mathrm{phase1}}(s_{i-1}, s_i, s_{i+1}) := (\gamma_i, q', 0).
    \]
  \item If the centre cell carries a left-moving mark, $q_i^\# = (q',L)$, then
    the head is leaving site $i$ to the left:
    \[
    U_{\mathrm{phase1}}(s_{i-1}, s_i, s_{i+1}) := (\gamma_i, \bot, 0).
    \]
  \item If the \emph{right neighbour} carries a left-moving mark and the centre
    cell is headless, $q_{i+1}^\# = (q',L)$ and $q_i^\# = \bot$, then the head
    arrives from the right:
    \[
    U_{\mathrm{phase1}}(s_{i-1}, s_i, s_{i+1}) := (\gamma_i, q', 0).
    \]
  \item In all remaining cases (cells that neither carry nor receive a head),
    the cell preserves its tape symbol and any non-directional head component
    and simply flips the phase bit $1 \mapsto 0$. Under the invariant stated
    below, this means
    \[
    U_{\mathrm{phase1}}(s_{i-1}, s_i, s_{i+1}) := (\gamma_i, q_i^*, 0)
    \]
    where $q_i^*$ is either $\bot$ or $q \in Q$ (never a directional mark).
\end{itemize}

The full local rule $U$ is obtained by combining the two phases into a single
radius-$1$ rule on $S^{\{-1,0,1\}}$, with the phase bit $p$ in the centre cell
determining which branch, $U_{\mathrm{phase0}}$ or $U_{\mathrm{phase1}}$,
applies at that cell.

\subsubsection{Encoding, decoding, and correctness}

We now define the WCC machine $W_M$ associated to $M$.

\begin{definition}[Associated WCC machine $W_M$]
\label{def:wm}
Let $M$ be a deterministic Turing machine as above. Define $W_M$ to be the
discrete WCC machine with:
\begin{itemize}
  \item dimension $d = 1$,
  \item index sets $\Lambda_n = \{-L(n),\dots,L(n)\}$ with $L(n)$ a fixed
    polynomial upper bound on $M$'s tape usage on inputs of length $n$,
  \item state set $S$ as in Definition~\ref{def:local_state} and radius $r = 1$,
    with neighbourhoods $N_1(i) = \{i-1,i,i+1\}$,
  \item local update rule $U : S^{\{-1,0,1\}} \to S$ as defined above,
  \item encoding $\mathrm{enc}_M : \Sigma^* \to S^{\Lambda_n}$ given by
    \[
    \mathrm{enc}_M(x)_i =
    \begin{cases}
      (x_1, q_0, 0), & i = 0,\\[2pt]
      (x_{i+1}, \bot, 0), & 1 \le i \le n-1,\\[2pt]
      (\blank, \bot, 0), & \text{otherwise},
    \end{cases}
    \]
    i.e.~the input word $x = x_1 \dots x_n$ is written on cells $0,\dots,n-1$,
    the head starts at $i=0$ in state $q_0$, and the rest of the tape is blank,
  \item decoding predicate $\mathrm{dec}_M : S^{\Lambda_n} \to \{\mathsf{run},0,1\}$ defined
    by
    \[
    \mathrm{dec}_M(C) =
    \begin{cases}
      1, & \text{if the simulated head is in }q_{\mathrm{acc}},\\
      0, & \text{if the simulated head is in }q_{\mathrm{rej}},\\
      \mathsf{run}, & \text{otherwise}.
    \end{cases}
    \]
\end{itemize}
\end{definition}

We first state an invariant that characterizes valid configurations in the
simulation.

\begin{lemma}[Simulation invariant]
\label{lem:invariant}
For the WCC machine $W_M$ on input $x \in \Sigma^n$, the following holds for all
$t \ge 0$:
\begin{itemize}
  \item At every \emph{even} time $2t$, all cells have phase bit $p=0$ and there
    is a unique head cell $h_t \in \Lambda_n$ such that
    $C_{2t}(h_t) = (\gamma_t(h_t), q_t, 0)$ with $q_t \in Q$, while all other
    cells have $q^\# = \bot$.
  \item At every \emph{odd} time $2t+1$, all cells have phase bit $p=1$ and
    there is a unique directional mark at some site $j_t \in \Lambda_n$ with
    $C_{2t+1}(j_t) = (\widetilde\gamma(j_t), (q',d), 1)$ for some $q' \in Q$
    and $d \in \{L,R\}$, while all other cells have
    $q^\# \in Q \cup \{\bot\}$ but never a directional mark.
\end{itemize}
\end{lemma}

\begin{proof}
The base case follows from the definition of $\mathrm{enc}_M$ at $t=0$. The step
from $2t$ to $2t+1$ uses $U_{\mathrm{phase0}}$, which can only introduce a
single directional mark at the unique head cell and flips all phases to $1$. The
step from $2t+1$ to $2t+2$ uses $U_{\mathrm{phase1}}$, which ensures that any
directional mark at the centre cell is cleared and exactly one neighbour
installs the corresponding non-directional head state; all other cells remain
without directional marks and all phases flip back to $0$. The uniqueness of the
marked site follows from the fact that there is exactly one directional mark at
odd times and the local rule never creates additional marks; a straightforward
case analysis on the neighbourhoods shows that no two cells can simultaneously
install a head state from the same marked neighbour. This preserves the
invariant.
\end{proof}

We now establish configuration correspondence and runtime.

\begin{lemma}[One Turing step equals two WCC steps]
\label{lem:one_step}
Let $(\gamma_t, h_t, q_t)$ be the Turing configuration of $M$ after $t$ steps
on input $x$, and let $C_{2t}$ be the configuration of $W_M$ after $2t$ update
steps with initial condition $C_0 = \mathrm{enc}_M(x)$. Then, for all $t \ge 0$,
the configuration $C_{2t}$ encodes $(\gamma_t, h_t, q_t)$ in the sense that:
\begin{itemize}
  \item For all $i \in \Lambda_n$, the tape symbol in $C_{2t}(i)$ equals
    $\gamma_t(i)$.
  \item There exists a unique $i = h_t$ such that the head component of
    $C_{2t}(i)$ equals $q_t$ and all other cells carry $\bot$ in the head
    component.
  \item All cells have phase bit $p=0$.
\end{itemize}
Moreover, applying two further WCC update steps transforms $C_{2t}$ into
$C_{2(t+1)}$ which encodes $(\gamma_{t+1}, h_{t+1}, q_{t+1})$.
\end{lemma}

\begin{proof}
The claim follows by induction on $t$. For $t=0$ the statement holds by
construction of $\mathrm{enc}_M$ and the invariant of Lemma~\ref{lem:invariant}.
Suppose the claim holds for some $t\ge 0$. In configuration $C_{2t}$ there is
exactly one cell $h_t$ with head state $q_t$ at phase $p=0$, and the tape
symbols agree with $\gamma_t(i)$ for all $i$.

The phase-$0$ part of the local rule applies $\delta(q_t,\gamma_t(h_t)) =
(q',\gamma',d)$ at cell $h_t$, overwriting the tape symbol with $\gamma'$ and
marking the head component as $(q',d)$ while flipping the phase to $1$. All
other cells preserve their tape symbols and have no head, and flip their phase
to $1$. After this first update, we obtain $C_{2t+1}$, in which the tape symbols
coincide with those of $\gamma_t$ except at $h_t$, where they match the updated
symbol $\gamma'$, and the unique directional head mark $(q',d)$ is at cell
$h_t$, in accordance with Lemma~\ref{lem:invariant}.

In the next update step (phase $1$), the radius-$1$ rule $U_{\mathrm{phase1}}$
moves the directional head mark one cell in direction $d$ to obtain a unique
head state $q'$ at either $h_t+1$ (if $d=R$) or $h_t-1$ (if $d=L$), clears the
directional mark, and restores all phase bits to $0$, producing $C_{2(t+1)}$.
The tape symbols are unchanged in this step. By definition of the Turing
transition, the new Turing configuration is $(\gamma_{t+1}, h_{t+1}, q_{t+1})$
with $\gamma_{t+1}$ and $(h_{t+1},q_{t+1})$ matching the updated symbol and
head location/state. Thus $C_{2(t+1)}$ encodes $(\gamma_{t+1}, h_{t+1},
q_{t+1})$ as claimed. This completes the induction.
\end{proof}

\begin{lemma}[Runtime and space overhead]
\label{lem:runtime}
Let $T_M(n)$ be the worst-case runtime of $M$ on inputs of length $n$. Then the
associated WCC machine $W_M$ satisfies:
\begin{itemize}
  \item For all $x \in \Sigma^n$, $W_M$ reaches an accepting configuration if
    and only if $M$ enters $q_{\mathrm{acc}}$ on $x$.
  \item The number of WCC update steps until halting satisfies
    \[
    T_{W_M}(n) = 2 T_M(n).
    \]
  \item The number of cells satisfies $|\Lambda_n| \le 2L(n)+1 \le
    \tilde{p}(n)$ for some polynomial $\tilde{p}$.
\end{itemize}
\end{lemma}

\begin{proof}
The first two claims follow directly from Lemma~\ref{lem:one_step} and the
definition of $\mathrm{dec}_M$: after $T_M(n)$ Turing steps, $M$ halts in either
$q_{\mathrm{acc}}$ or $q_{\mathrm{rej}}$, and the corresponding encoded
configuration appears at WCC time $2T_M(n)$ with phase $p=0$; the decoding
predicate tests for the presence of $q_{\mathrm{acc}}$ and thus detects
acceptance exactly when $M$ accepts.

For the space bound, by hypothesis $M$ uses at most $L(n)$ nonblank tape cells
on inputs of length $n$, so choosing $\Lambda_n = \{-L(n),\dots,L(n)\}$
suffices. Since $L(n)$ is bounded by a polynomial in $n$, we have
$|\Lambda_n| = 2L(n)+1 \le \tilde{p}(n)$ for a polynomial $\tilde{p}$. Each
update step applies the local rule $U$ independently at each cell, so a single
global update can be carried out with work $O(|\Lambda_n|) =
\mathrm{poly}(n)$, which respects the polynomial-time cost per pass assumed in
Definition~\ref{def:pwcc}.
\end{proof}

We can now state the main equivalence theorem for the simulation direction
TM $\Rightarrow$ WCC.

\begin{theorem}[Simulation of polynomial-time Turing machines by WCC]
\label{thm:tm_to_wcc}
Let $L \subseteq \Sigma^*$ be a language decided by a deterministic Turing
machine $M$ in time $T_M(n) \le c n^k$ for some constants $c,k \ge 1$. Then the
associated WCC machine $W_M$ defined in Definition~\ref{def:wm} decides $L$ in
time
\[
T_{W_M}(n) = 2 T_M(n) \le 2 c n^k
\]
on inputs of length $n$. In particular,
\[
P_{\mathrm{Turing}} \subseteq P_{\mathrm{WCC}}.
\]
\end{theorem}

\begin{proof}
By Lemma~\ref{lem:runtime}, $W_M$ accepts exactly the inputs accepted by $M$ and
does so in at most $2 T_M(n)$ global update steps on inputs of length $n$. The
bound $T_M(n) \le c n^k$ implies $T_{W_M}(n) \le 2 c n^k$, so $W_M$ is a
polynomial-time WCC machine in the sense of Definition~\ref{def:pwcc}. Therefore
$L \in P_{\mathrm{WCC}}$. Since $L$ was arbitrary in $P_{\mathrm{Turing}}$, we
obtain $P_{\mathrm{Turing}} \subseteq P_{\mathrm{WCC}}$.
\end{proof}

\subsection{Simulation of WCC by Turing machines}

We now recall the converse direction: any discrete WCC machine can be simulated
by a deterministic Turing machine with polynomial overhead. This is essentially
a standard ``Turing simulation of a cellular automaton'' argument.

\begin{theorem}[Simulation of WCC by Turing machines]
\label{thm:wcc_to_tm}
Let $W$ be a discrete WCC machine that decides a language $L$ in polynomial
time as in Definition~\ref{def:pwcc}, with runtime bound $T_W(n) \le n^k$ and
$|\Lambda_n| \le p(n)$ for some polynomials $k,p$. Then there exists a
deterministic Turing machine $M_W$ that decides $L$ in time polynomially
bounded in $n$. In particular,
\[
P_{\mathrm{WCC}} \subseteq P_{\mathrm{Turing}}.
\]
\end{theorem}

\begin{proof}
We sketch the standard simulation; details can be made fully explicit if
desired.

Fix $n$ and the corresponding index set $\Lambda_n$. A configuration
$C_t : \Lambda_n \to S$ of $W$ can be encoded on the tape of a Turing machine as
a finite word over $\Gamma_{\mathrm{enc}} := S \cup \{\#\}$ by listing the cell
states in some canonical order $i_1,\dots,i_{|\Lambda_n|}$ (e.g.~lexicographic
order on $\Lambda_n$) separated by a delimiter $\#$. This yields an encoding
$\mathrm{code}(C_t)$ of length $O(|\Lambda_n|)$.

To simulate one global WCC update $C_t \mapsto C_{t+1}$, the Turing machine
$M_W$ maintains the current configuration encoding on its work tape and proceeds
in two phases:
\begin{enumerate}
  \item It copies the current configuration to an auxiliary work area. This
    takes time $O(|\Lambda_n|)$.
  \item It enumerates each cell, uses the uniform lattice/neighborhood
    algorithms to retrieve its at most $(2r+1)^d$ local symbols,
    computes the fixed rule $U$, and writes the next state. The local
    lookup and update have polynomial cost in $n$; a simple tape
    implementation may take $O(|\Lambda_n|^2\,\mathrm{poly}(n))$
    time per round. No constant-time random access is assumed.
\end{enumerate}
By assumption $|\Lambda_n| \le p(n)$ and one global update can be effected with
cost polynomial in $|\Lambda_n|$, i.e.~$\mathrm{poly}(p(n))$. Repeating this
procedure for $T_W(n) \le n^k$ steps yields an overall Turing runtime
\[
T_{M_W}(n) \le T_W(n) \cdot \mathrm{poly}(|\Lambda_n|)
\le n^k \cdot \mathrm{poly}(p(n))
= \mathrm{poly}(n).
\]
The Turing machine $M_W$ accepts exactly when the simulated configuration
satisfies $\mathrm{dec}(C_t) = 1$, which can be checked during or after each
simulated global update by scanning the encoding for an accepting pattern. Thus
$M_W$ decides the same language $L$ in polynomial time, and $L \in
P_{\mathrm{Turing}}$. Since $L$ was arbitrary in $P_{\mathrm{WCC}}$, it follows
that $P_{\mathrm{WCC}} \subseteq P_{\mathrm{Turing}}$.
\end{proof}

\subsection{Equivalence of $P_{\mathrm{WCC}}$ and $P_{\mathrm{Turing}}$}

Combining Theorem~\ref{thm:tm_to_wcc} and Theorem~\ref{thm:wcc_to_tm} yields
the desired identification.

\begin{corollary}[Polynomial-time equivalence]
\label{cor:pwcc_equals_pturing}
Under the discrete WCC machine model of Definition~\ref{def:wcc_machine} and the
standard deterministic Turing machine model of Definition~\ref{def:tm}, the
polynomial-time classes coincide:
\[
P_{\mathrm{WCC}} = P_{\mathrm{Turing}}.
\]
\end{corollary}

% WIN-11: physical resource equivalence is separate from discrete equivalence.
\begin{remark}[Physical WCT realization is a separate obligation]
\label{rem:wct_realization}
The equality $P_{\mathrm{WCC}}^{\mathrm{disc}}=P$ concerns only the
uniform finite-state abstraction of Definition~\ref{def:wcc_machine}.
It does \emph{not} imply $P_{\mathrm{WCC}}^{\mathrm{phys}}=P$.
To deduce that equality for a curvature-constrained WCT device family,
one must additionally construct, uniformly in input length and
for every required local rule, a stable encoding $s\mapsto\psi_s$ and
a physical evolution satisfying:
\begin{enumerate}
\item the exact code-state transition compatibility of
Theorem~\ref{thm:wave-symbol-isomorphism};
\item polynomially bounded preparation, run time, volume,
energy/curvature expenditure, precision and readout costs;
\item correct uniform simulation of all polynomial-time discrete
WCC computations, including their verifier machines.
\end{enumerate}
These are additional model-specific hypotheses. They are neither
consequences of locality and finite precision nor established by the
Turing simulations proved in this section. The reverse inclusion
$P_{\mathrm{WCC}}^{\mathrm{phys}}\subseteq P$ would likewise require
a precise representation and polynomial-time classical simulation of
all admissible physical operations.


\end{remark}

% ============================
% WCC verifiers and NP equivalence
% ============================

\section{WCC Verifiers and NP Equivalence}
\label{sec:np_wcc_equiv}

In this section we lift the polynomial-time equivalence of
Corollary~\ref{cor:pwcc_equals_pturing} from decision procedures to
verifier-based nondeterministic computation. We define an NP-type class
$NP_{\mathrm{WCC}}$ using polynomial-time WCC verifiers and show that, under the
same discrete model assumptions, it coincides with the classical
nondeterministic polynomial-time class $NP_{\mathrm{Turing}}$.

\subsection{Classical NP via Turing verifiers}

We first recall the standard verifier-based definition of $NP$.

Let $\Sigma$ be a finite input alphabet and let $\Gamma_{\mathrm{wit}}$ be a
finite witness alphabet. For $x \in \Sigma^n$ and
$y \in \Gamma_{\mathrm{wit}}^{m}$ we write $(x,y)$ for the pair instance, and
$N := n + m$ for the total input length.

\begin{definition}[Polynomial-time Turing verifier]
\label{def:tm_verifier}
A \emph{polynomial-time Turing verifier} for a language $L \subseteq \Sigma^*$
is a deterministic Turing machine
\[
V_{\mathrm{T}} = (Q, \Gamma, \Sigma \cup \Gamma_{\mathrm{wit}}, \delta, q_0,
q_{\mathrm{acc}}, q_{\mathrm{rej}})
\]
such that:
\begin{itemize}
  \item $V_{\mathrm{T}}$ takes as input a pair $(x,y)$ with $x \in \Sigma^n$,
    $y \in \Gamma_{\mathrm{wit}}^{m}$, encoded in some fixed, efficiently
    decodable format on its tape.
  \item There exist polynomials $p,q$ such that for all $x \in \Sigma^n$ and
    all $y$,
    \[
    T_{V_{\mathrm{T}}}(|x|+|y|) \le p(|x|+|y|)
    \]
    and
    \[
    x \in L
    \iff
    \exists y \in \Gamma_{\mathrm{wit}}^{m},\ m \le q(|x|) \text{ such that }
    V_{\mathrm{T}}(x,y) \text{ accepts.}
    \]
\end{itemize}
\end{definition}

\begin{definition}[Classical nondeterministic polynomial time]
\label{def:npturing}
A language $L \subseteq \Sigma^*$ is in $NP_{\mathrm{Turing}}$ if there exists a
polynomial-time Turing verifier $V_{\mathrm{T}}$ for $L$ in the sense of
Definition~\ref{def:tm_verifier}. The class $NP_{\mathrm{Turing}}$ is the set of
all such languages.
\end{definition}

This definition is equivalent to the usual one via nondeterministic Turing
machines, but we will not need that equivalence here.

\subsection{WCC verifiers and $NP_{\mathrm{WCC}}$}

We now define polynomial-time WCC verifiers and the associated NP-type class.

\begin{definition}[Polynomial-time WCC verifier]
\label{def:wcc_verifier}
Let $\Sigma$ be a finite input alphabet and $\Gamma_{\mathrm{wit}}$ a finite
witness alphabet. A \emph{polynomial-time WCC verifier} for a language
$L \subseteq \Sigma^*$ is a discrete WCC machine
\[
V_{\mathrm{W}} = (\{\Lambda_N\}_{N\ge 1}, S, r, U, \mathrm{enc}_{\mathrm{W}},
\mathrm{dec}_{\mathrm{W}})
\]
in the sense of Definition~\ref{def:wcc_machine}, together with polynomials
$p,q$ such that:
\begin{itemize}
  \item For each instance $x \in \Sigma^n$ and candidate witness
    $y \in \Gamma_{\mathrm{wit}}^{m}$ with $m \le q(n)$, the pair $(x,y)$ is
    encoded as an initial configuration
    \[
    C_0 = \mathrm{enc}_{\mathrm{W}}(x,y) \in S^{\Lambda_N},
    \qquad N := n + m,
    \]
    for some index set $\Lambda_N$ with $|\Lambda_N| \le p(N)$.
  \item The WCC runtime $T_{V_{\mathrm{W}}}(N)$, defined as the worst-case
    number of global update steps until $\mathrm{dec}_{\mathrm{W}}(C_t)$
    first returns $0$ or $1$ on inputs of total length $N$, satisfies
    \[
    T_{V_{\mathrm{W}}}(N) \le p(N)
    \]
    for all $N$.
  \item For all $x \in \Sigma^n$,
    \[
    x \in L
    \iff
    \exists y \in \Gamma_{\mathrm{wit}}^{m},\ m \le q(n) \text{ such that }
    V_{\mathrm{W}}(x,y) \text{ accepts,}
    \]
    where acceptance means $\mathrm{dec}_{\mathrm{W}}(C_t) = 1$ for some
    $t \le T_{V_{\mathrm{W}}}(|x|+|y|)$.
\end{itemize}
\end{definition}

\begin{definition}[WCC nondeterministic polynomial time]
\label{def:npwcc}
A language $L \subseteq \Sigma^*$ is in $NP_{\mathrm{WCC}}$ if there exists a
polynomial-time WCC verifier $V_{\mathrm{W}}$ for $L$ in the sense of
Definition~\ref{def:wcc_verifier}. The class $NP_{\mathrm{WCC}}$ is the set of
all such languages.
\end{definition}

By construction, $NP_{\mathrm{WCC}}$ is defined using exactly the same
witness-length and runtime scaling as $NP_{\mathrm{Turing}}$, but with WCC
machines in place of Turing machines.

\subsection{Equivalence of $NP_{\mathrm{WCC}}$ and $NP_{\mathrm{Turing}}$}

We now prove that, under the discrete WCC model, the verifier-based NP class
coincides with the classical one.

\begin{theorem}[NP equivalence]
\label{thm:np_equiv}
Under the discrete WCC machine model of Definition~\ref{def:wcc_machine} and the
Turing verifier model of Definition~\ref{def:tm_verifier}, the
nondeterministic polynomial-time classes coincide:
\[
NP_{\mathrm{WCC}} = NP_{\mathrm{Turing}}.
\]
\end{theorem}

\begin{proof}
We prove both inclusions.

\paragraph{Inclusion $NP_{\mathrm{WCC}} \subseteq NP_{\mathrm{Turing}}$.}

Suppose $L \in NP_{\mathrm{WCC}}$. Then there exists a polynomial-time WCC
verifier $V_{\mathrm{W}}$ for $L$ with polynomials $p,q$ as in
Definition~\ref{def:wcc_verifier}. For each pair $(x,y)$ with $x \in \Sigma^n$,
$y \in \Gamma_{\mathrm{wit}}^{m}$, $m \le q(n)$, the machine $V_{\mathrm{W}}$
runs in at most $p(N)$ global update steps, where $N := n+m$, on a lattice
$\Lambda_N$ with $|\Lambda_N| \le p(N)$.

By Theorem~\ref{thm:wcc_to_tm}, there exists a deterministic Turing machine
$M_{V_{\mathrm{W}}}$ that simulates $V_{\mathrm{W}}$ on inputs of total length
$N$ in time bounded by
\[
T_{M_{V_{\mathrm{W}}}}(N) \le p(N) \cdot \mathrm{poly}(|\Lambda_N|)
\le p(N) \cdot \mathrm{poly}(p(N))
= \mathrm{poly}(N).
\]
Define a Turing verifier $V_{\mathrm{T}}$ that, on input $(x,y)$, first
constructs the encoding $\mathrm{enc}_{\mathrm{W}}(x,y)$ on its tape and then
simulates $M_{V_{\mathrm{W}}}$ for $p(N) \cdot \mathrm{poly}(p(N))$ steps,
accepting if and only if $V_{\mathrm{W}}(x,y)$ accepts. This adds at most an
additional polynomial overhead in $N$ for the encoding step, so $V_{\mathrm{T}}$
runs in time polynomial in $N$.

By construction, for either verifier $V$ in
$\{V_{\mathrm{W}},V_{\mathrm{T}}\}$ the same predicate is decided:
\[
 x\in L\quad\Longleftrightarrow\quad
 \exists y\ \bigl(|y|\le q(|x|)\ \wedge\ V(x,y)=1\bigr).
\]
Thus $V_{\mathrm{T}}$ is a polynomial-time Turing verifier for $L$ in the sense
of Definition~\ref{def:tm_verifier}, and hence $L \in NP_{\mathrm{Turing}}$.
Since $L$ was arbitrary in $NP_{\mathrm{WCC}}$, we obtain
$NP_{\mathrm{WCC}} \subseteq NP_{\mathrm{Turing}}$.

\paragraph{Inclusion $NP_{\mathrm{Turing}} \subseteq NP_{\mathrm{WCC}}$.}

Conversely, suppose $L \in NP_{\mathrm{Turing}}$. Then there exists a
polynomial-time Turing verifier $V_{\mathrm{T}}$ for $L$ with polynomials $p,q$
as in Definition~\ref{def:tm_verifier}. For each pair $(x,y)$ with
$x \in \Sigma^n$, $y \in \Gamma_{\mathrm{wit}}^{m}$, $m \le q(n)$, the machine
$V_{\mathrm{T}}$ halts in at most $p(N)$ steps with $N := n+m$.

Apply the construction of Theorem~\ref{thm:tm_to_wcc} to $V_{\mathrm{T}}$,
treating $(x,y)$ as a single input string of length $N$. This yields a discrete
WCC machine $W_{V_{\mathrm{T}}}$ that simulates $V_{\mathrm{T}}$ with at most a
constant-factor overhead in time, i.e.
\[
T_{W_{V_{\mathrm{T}}}}(N) \le c\, T_{V_{\mathrm{T}}}(N) \le c\, p(N)
\]
for some constant $c\ge 1$, on index sets $\Lambda_N$ with $|\Lambda_N| \le
\tilde{p}(N)$ for some polynomial $\tilde{p}$. We define the encoding map
$\mathrm{enc}_{\mathrm{W}}$ for $W_{V_{\mathrm{T}}}$ by viewing $(x,y)$ simply
as a word over the combined alphabet $\Sigma \cup \Gamma_{\mathrm{wit}}$ and
placing it on the initial tape segment as in Definition~\ref{def:wm}. The
decoding predicate $\mathrm{dec}_{\mathrm{W}}$ tests for the presence of the
accepting state $q_{\mathrm{acc}}$ in the head component, exactly as in
Definition~\ref{def:wm}.

By Lemma~\ref{lem:runtime} and Theorem~\ref{thm:tm_to_wcc}, we then have
\[
(x,y) \text{ accepted by } V_{\mathrm{T}}
\iff
(x,y) \text{ accepted by } W_{V_{\mathrm{T}}}
\]
for all pairs $(x,y)$. Thus, with $V$ equal to either
$V_{\mathrm{T}}$ or $W_{V_{\mathrm{T}}}$,
\[
 x\in L\quad\Longleftrightarrow\quad
 \exists y\ \bigl(|y|\le q(|x|)\ \wedge\ V(x,y)=1\bigr).
\]
Thus $W_{V_{\mathrm{T}}}$ is a polynomial-time WCC verifier for $L$ with
witness bound $q(\cdot)$ and runtime bounded by a polynomial in $N$, and hence
$L \in NP_{\mathrm{WCC}}$.

Since $L$ was arbitrary in $NP_{\mathrm{Turing}}$, we obtain
$NP_{\mathrm{Turing}} \subseteq NP_{\mathrm{WCC}}$.

\paragraph{Conclusion.}

We have shown $NP_{\mathrm{WCC}} \subseteq NP_{\mathrm{Turing}}$ and
$NP_{\mathrm{Turing}} \subseteq NP_{\mathrm{WCC}}$, hence
\[
NP_{\mathrm{WCC}} = NP_{\mathrm{Turing}}.
\]
This completes the proof.
\end{proof}




% ============================
% Curvature-bounded channel capacity in WCT
% ============================

\section{Curvature-Bounded Channel Capacity in WCT}
\label{sec:curvature_capacity}

In this section we make precise the intuition that a wave-confinement device with finite volume, energy budget, and runtime cannot sustain arbitrarily many independent information ``lanes'' or arbitrarily fast switching between them. Under the rest-energy curvature identity and the dimensional lock on timescales, we obtain a polynomial upper bound on the spatiotemporal information-processing capacity of a WCT cavity.

\subsection{Geometric and energetic setup}

Let $\Gamma$ be a smooth closed characteristic loop of a confined field configuration with arclength parameter $s \in [0,L)$, total length
\[
L = \oint_\Gamma ds,
\]
and Frenet--Serret curvature $\kappa(s)$ and torsion $\tau(s)$. We define the local curvature-rate scalar
\begin{equation}
  \sigma(s)^2 := \kappa(s)^2 + \tau(s)^2.
\end{equation}

Let $w(s) > 0$ be a weight function induced, for example, by an energy density or probability density derived from $\abs{\psi}^2$ along $\Gamma$. We write the corresponding weighted average of a quantity $f$ as
\begin{equation}
  \label{eq:weighted_average}
  \langle f \rangle_w
  := \frac{\displaystyle \oint_\Gamma w(s)\,f(s)\,ds}{\displaystyle \oint_\Gamma w(s)\,ds}.
\end{equation}

\begin{definition}[Rest energy / mass from curvature]
\label{def:rest_energy_curvature}
We say that a confined field configuration obeys the \emph{rest energy curvature identity} if its rest energy and inertial mass are given by
\begin{equation}
  \label{eq:rest_energy_identity}
  E_{\mathrm{rest}} = \hbar c\,\langle \sigma \rangle_w,
  \qquad
  m = \frac{\hbar}{c}\,\langle \sigma \rangle_w,
\end{equation}
for some positive weight $w$ along closed characteristic loops of the configuration.
\end{definition}

In a bulk description, we consider a bounded cavity $\Omega \subset \mathbb{R}^3$ with volume $V := \abs{\Omega}$ confining a complex scalar field
\[
\psi : \Omega \times \mathbb{R} \to \mathbb{C},
\]
which we assume satisfies a WCT-type curvature-locking PDE. We denote by $E_{\mathrm{rest}}(\psi)$ the total rest energy of $\psi$ in $\Omega$ and impose a resource bound
\begin{equation}
  \label{eq:energy_budget}
  E_{\mathrm{rest}}(\psi) \le E_{\max}.
\end{equation}

\subsection{Curvature, bandlimit, and spatial capacity}

We first connect the curvature/energy budget to an effective bandlimit on the field.

\begin{definition}[Effective wavenumber]
\label{def:effective_k}
Let $\widehat{\psi}(\mathbf{k})$ denote the spatial Fourier transform of $\psi(\cdot,t)$ at a fixed time. We say that $\psi$ is \emph{effectively band-limited} to radius $k_{\max}$ if
\begin{equation}
  \widehat{\psi}(\mathbf{k}) = 0
  \quad\text{whenever}\quad
  \norm{\mathbf{k}} > k_{\max}.
\end{equation}
We denote by $k_{\mathrm{eff}}$ an effective wavenumber scale associated to $\psi$, for example a suitable norm or average of $\mathbf{k}$ weighted by $\abs{\widehat{\psi}}^2$.
\end{definition}

\begin{definition}[Curvature--wavenumber coupling]
\label{def:curvature_k_coupling}
We assume that for confined WCT configurations the effective wavenumber is controlled by the curvature-rate average in the sense that there exists a constant $C_1 > 0$ with
\begin{equation}
  \label{eq:curvature_k_coupling}
  k_{\mathrm{eff}} \le C_1 \,\langle \sigma \rangle_w.
\end{equation}
\end{definition}

Combining the energy identity with the assumed coupling bounds the
\emph{effective} wavenumber, but does not by itself bound the maximum
Fourier frequency or imply compact spectral support. A separate sharp
bandwidth hypothesis is needed for finite-mode counting.

\begin{lemma}[Effective wavenumber estimate; conditional sharp cutoff]
\label{lem:bandlimit}
Assume the rest-energy identity \eqref{eq:rest_energy_identity}, the
coupling \eqref{eq:curvature_k_coupling}, and
$E_{\mathrm{rest}}(\psi)\le E_{\max}$. Then
\begin{equation}
 k_{\mathrm{eff}}\le C_1\frac{E_{\max}}{\hbar c}.
 \label{eq:effective_k_energy}
\end{equation}
If one \emph{additionally assumes} a uniform spectral-support bound
$\operatorname{supp}\widehat\psi\subseteq\{\|k\|\le
k_{\max}\}$ and the stronger coupling
$k_{\max}\le C_1\langle\sigma\rangle_w$, then
\begin{equation}
 k_{\max}\le C_1\frac{E_{\max}}{\hbar c}.
 \label{eq:kmax_energy}
\end{equation}
The second conclusion is conditional on the sharp-support premise;
it does not follow from the effective-wavenumber estimate alone.
\end{lemma}



\begin{proof}
The energy identity and bound give
$\langle\sigma\rangle_w\le E_{\max}/(\hbar c)$.
Substitution into the corresponding assumed coupling yields
\eqref{eq:effective_k_energy}; using the \emph{additional}
sharp-cutoff coupling yields \eqref{eq:kmax_energy}.
\end{proof}



We now translate the bandlimit into a bound on the number of independent spatial degrees of freedom (``lanes'').

\begin{lemma}[Curvature-bounded spatial channel capacity]
\label{lem:spatial_capacity}
Let $\Omega$ be a fixed periodic cubical cavity with volume $V$, and
let the computation use only its Fourier modes satisfying
$\norm{\mathbf{k}}\le k_{\max}$. Assume $k_{\max}$ is at
least the fundamental nonzero wavenumber of the cavity. Then there
is a geometric constant $C_2>0$ such that the number
$N_{\mathrm{lanes}}$ of these modes obeys
\begin{equation}
  \label{eq:spatial_capacity}
  N_{\mathrm{lanes}} \;\le\; C_2\,V\,k_{\max}^3.
\end{equation}
In particular, under the energy budget $E_{\mathrm{rest}}(\psi) \le E_{\max}$ we have
\begin{equation}
  \label{eq:spatial_capacity_energy}
  N_{\mathrm{lanes}} \;\le\; C_2\,V\,
  \Big( C_1 \frac{E_{\max}}{\hbar c} \Big)^3.
\end{equation}
\end{lemma}

\begin{proof}
For a periodic cubic cavity, allowed wavenumbers form a lattice of
spacing proportional to $V^{-1/3}$. Counting the lattice points
inside the spectral ball gives, for $k_{\max}$ at least the
fundamental wavenumber, an upper bound proportional to $Vk_{\max}^3$.
The leading-volume approximation is
\[
N_{\mathrm{lanes}} \sim \frac{V}{(2\pi)^3} \,\frac{4\pi}{3}k_{\max}^3
= C_2\,V\,k_{\max}^3,
\]
for an appropriate constant $C_2 > 0$. Substituting the bound \eqref{eq:kmax_energy} from Lemma~\ref{lem:bandlimit} gives \eqref{eq:spatial_capacity_energy}.
\end{proof}

Here $N_{\mathrm{lanes}}$ counts orthogonal spatial Fourier modes,
not all metastable states in a nonlinear cavity. Its identification with
independently addressable code lanes is an additional modeling
hypothesis, not a general consequence of the mode count.

\subsection{Temporal dimensional lock and spatiotemporal capacity}

We now incorporate the temporal side of the dimensional lock. The key assumption is that curvature-locked code states cannot be switched arbitrarily fast.

\begin{definition}[Dimensional lock in time]
\label{def:temporal_lock}
Let $\Delta t_{\min}$ denote the minimal dwell time required for a curvature-locked configuration to relax from one metastable code state to a distinct one while remaining within the WCT code manifold. We assume there exists a constant $C_3 > 0$ such that
\begin{equation}
  \label{eq:temporal_lock}
  \Delta t_{\min} \;\ge\; \frac{C_3}{\omega_{\max}},
\end{equation}
where $\omega_{\max}$ is a characteristic maximal angular frequency of the band-limited field. For propagating modes we assume a dispersion relation of the form
\begin{equation}
  \label{eq:dispersion}
  \omega_{\max} \;\le\; c\,k_{\max},
\end{equation}
with $c$ the relevant propagation speed.
\end{definition}

\begin{lemma}[Per-lane temporal update bound]
\label{lem:lane_update_bound}
Under the assumptions of Lemma~\ref{lem:spatial_capacity} and Definition~\ref{def:temporal_lock}, over any time window of duration $T > 0$ a single lane can undergo at most
\begin{equation}
  \label{eq:steps_per_lane}
  N_{\mathrm{steps\ per\ lane}} \;\le\; \frac{T}{\Delta t_{\min}}
  \;\le\; \frac{T\,c\,k_{\max}}{C_3}
\end{equation}
distinct metastable state transitions.
\end{lemma}

\begin{proof}
The first inequality is the definition of $N_{\mathrm{steps\ per\ lane}}$. The second follows by combining \eqref{eq:temporal_lock} and \eqref{eq:dispersion}:
\[
\Delta t_{\min} \ge \frac{C_3}{\omega_{\max}}
\ge \frac{C_3}{c\,k_{\max}}
\quad\Longrightarrow\quad
\frac{T}{\Delta t_{\min}} \le \frac{T\,c\,k_{\max}}{C_3}.
\]
\end{proof}

Combining the spatial and temporal bounds yields a spatiotemporal capacity limit.

\begin{theorem}[Curvature-bounded spatiotemporal channel capacity]
\label{thm:spacetime_capacity}
Let $\Omega \subset \mathbb{R}^3$ be a bounded WCT cavity of volume $V$, confining a field $\psi$ with rest-energy bounded by $E_{\max}$, and suppose the sharp spectral-cutoff, periodic-cavity, independent-lane,
and minimum-dwell-time assumptions of
Lemmas~\ref{lem:bandlimit}, \ref{lem:spatial_capacity}, and
\ref{lem:lane_update_bound} hold. Then, over any time window of duration $T > 0$, the total number $N_{\mathrm{updates}}$ of distinct metastable lane transitions obeys
\begin{equation}
  \label{eq:Nupdates_kmax}
  N_{\mathrm{updates}} \;\le\; C_4\,V\,T\,k_{\max}^4,
\end{equation}
for some constant $C_4 > 0$, and hence
\begin{equation}
  \label{eq:Nupdates_energy}
  N_{\mathrm{updates}} \;\le\; K\,V\,T\,E_{\max}^4,
\end{equation}
for a constant $K > 0$ depending only on $C_1,C_2,C_3,\hbar,c$.
\end{theorem}

\begin{proof}
By Lemma~\ref{lem:spatial_capacity}, the number of lanes satisfies
\[
N_{\mathrm{lanes}} \le C_2\,V\,k_{\max}^3.
\]
By Lemma~\ref{lem:lane_update_bound}, each lane can switch at most $T c k_{\max}/C_3$ times in the interval of duration $T$. Thus the total number of lane updates obeys
\[
N_{\mathrm{updates}}
\;\le\; N_{\mathrm{lanes}} \cdot N_{\mathrm{steps\ per\ lane}}
\;\le\; \Big(C_2\,V\,k_{\max}^3\Big)\,\Big(\frac{T\,c\,k_{\max}}{C_3}\Big)
= C_4\,V\,T\,k_{\max}^4,
\]
where $C_4 := C_2 c / C_3$. Substituting the energy-dependent bound $k_{\max} \le C_1 E_{\max} / (\hbar c)$ from Lemma~\ref{lem:bandlimit} yields
\[
k_{\max}^4 \;\le\; C_1^4 \,\frac{E_{\max}^4}{\hbar^4 c^4},
\]
and hence
\[
N_{\mathrm{updates}} \;\le\; C_4\,V\,T\,k_{\max}^4
\;\le\; C_4\,V\,T\,C_1^4 \frac{E_{\max}^4}{\hbar^4 c^4}
\;=\; K\,V\,T\,E_{\max}^4
\]
for $K := C_4 C_1^4 / (\hbar^4 c^4)$.
\end{proof}

\begin{remark}[Polynomial resource scaling]
\label{rem:poly_resources}
Conditional on the spectral-support, mode-counting, and minimum
dwell-time assumptions, Theorem~\ref{thm:spacetime_capacity} shows that, for fixed physical resources $(V,T,E_{\max})$, the total number of resolvable metastable transitions in a WCT cavity is bounded by a polynomial in these parameters. In particular, if a family of WCT-based computing devices indexed by input size $n$ is constrained so that
\[
V(n),\; T(n),\; E_{\max}(n) \;\le\; \mathrm{poly}(n),
\]
then the total number $N_{\mathrm{updates}}(n)$ of elementary computational events across all lanes is also bounded by $\mathrm{poly}(n)$. Thus curvature confinement and dimensional lock enforce a physical bottleneck on the achievable computational complexity as a function of input size.
\end{remark}

















% ============================
% Physics-bounded computation and conditional NP safety
% ============================

\section{Physics-Bounded Computation and Conditional NP Safety}
\label{sec:phys_np_safety}

In this section we combine the discrete WCC--Turing equivalence with the
curvature/energy channel-capacity bound from the WCT model to obtain a
conditional statement about the computational power of any AGI implemented
on WCT hardware under the rails (L,P,V,C,S,R).

The next bound is a \emph{conditional physical hypothesis} derived
algebraically in the preceding section under the stated sharp-bandwidth,
mode-counting, and dwell-time premises; it has not been established
universally for arbitrary WCT PDE evolutions:

\begin{theorem}[Conditional channel-capacity step bound]
\label{thm:channel_capacity}
Let a WCT device occupy a cavity of volume $V > 0$ over a time interval
$[0,T]$ with a rest-energy ceiling $E_{\max} > 0$.
Assume the sharp spectral cutoff and the uniform mode-counting and
minimum-dwell-time hypotheses of Section~\ref{sec:curvature_capacity}.
Then the total number of distinct metastable lane updates
(individual local state transitions) $N_{\mathrm{updates}}$ that can occur
in $[0,T]$ obeys
\begin{equation}
  \label{eq:Nupdates_energy_restated}
  N_{\mathrm{updates}} \;\le\; K\,V\,T\,E_{\max}^4,
\end{equation}
for some constant $K > 0$ depending only on the WCT parameters
$(\hbar,c,\xi,\dots)$ and not on the particular device or algorithm.
\end{theorem}

We now formalize what it means for a family of physical devices to be
\emph{polynomially bounded} in WCT resources.

\begin{definition}[Polynomial-resource WCT device family]
\label{def:poly_resource_wct}
Let $(D_n)_{n\ge 1}$ be a family of WCT devices intended to process inputs
of length $n$. For each $n$, let $V(n)$ denote the device volume,
$T_{\mathrm{phys}}(n)$ the physical run time, and $E_{\max}(n)$ a rest-energy
ceiling during the computation.

We say that $(D_n)$ is a \emph{polynomial-resource WCT device family} if
there exist constants $C_V,C_T,C_E > 0$ and exponents $a,b,c \ge 0$ such that
for all $n$,
\begin{equation}
  \label{eq:poly_resource_scaling}
  V(n) \;\le\; C_V n^{a},\qquad
  T_{\mathrm{phys}}(n) \;\le\; C_T n^{b},\qquad
  E_{\max}(n) \;\le\; C_E n^{c}.
\end{equation}
\end{definition}

Combining Definition~\ref{def:poly_resource_wct} with
Theorem~\ref{thm:channel_capacity} yields a polynomial bound on the total
number of usable primitive updates.

\begin{proposition}[Physical step ceiling under WCT rails]
\label{prop:phys_step_ceiling}
Let $(D_n)_{n\ge 1}$ be a polynomial-resource WCT device family in the sense
of Definition~\ref{def:poly_resource_wct}. Then there exist constants
$C' > 0$ and $d \ge 0$ such that, for all $n$,
\begin{equation}
  \label{eq:Nupdates_poly}
  N_{\mathrm{updates}}(n) \;\le\; C' n^{d}.
\end{equation}
In particular, any discrete WCC machine $W$ faithfully implemented by such
a device family (with every nonhalting global round incurring at least one
counted elementary event) has worst-case runtime
$T_W(n)$ bounded by a polynomial in $n$.
\end{proposition}

\begin{proof}
By Definition~\ref{def:poly_resource_wct} we have
\[
V(n) \le C_V n^{a},\quad
T_{\mathrm{phys}}(n) \le C_T n^{b},\quad
E_{\max}(n) \le C_E n^{c}.
\]
Plugging these into the channel-capacity bound
\eqref{eq:Nupdates_energy} gives
\[
\begin{aligned}
N_{\mathrm{updates}}(n)
&\le K\,V(n)\,T_{\mathrm{phys}}(n)\,E_{\max}(n)^4\\
&\le K \cdot (C_V n^{a}) \cdot (C_T n^{b}) \cdot (C_E n^{c})^4\\
&= K C_V C_T C_E^4 \; n^{\,a + b + 4c}.
\end{aligned}
\]
Setting $C' := K C_V C_T C_E^4$ and $d := a + b + 4c$ yields
\eqref{eq:Nupdates_poly}. If every nonhalting global round incurs at least one counted
elementary event, then $T_W(n)\le N_{\mathrm{updates}}(n)+1$;
therefore $T_W$ is polynomially bounded. This counting premise must
hold for the particular implementation and clocking convention.
\end{proof}

We now combine the physical step ceiling with the discrete model equivalences
$P_{\mathrm{WCC}} = P_{\mathrm{Turing}}$ and
$NP_{\mathrm{WCC}} = NP_{\mathrm{Turing}}$ from
Sections~\ref{sec:tm_wcc_equiv}--\ref{sec:np_wcc_equiv} to obtain a
conditional NP-safety statement.

\begin{theorem}[Conditional NP-safety of WCT-bounded AGI]
\label{thm:conditional_np_safety}
Assume:
\begin{enumerate}[label=(\roman*)]
  \item The WCT rails (L,P,V,C,S,R) correctly describe the physically
        realizable dynamics of any WCT device, and the channel-capacity
        bound of Theorem~\ref{thm:channel_capacity} holds.
  \item All physically realizable discrete computation on such devices is
        faithfully modeled by discrete WCC machines in the sense of
        Definition~\ref{def:wcc_machine}, with equivalences
        \[
        P_{\mathrm{WCC}} = P_{\mathrm{Turing}},\qquad
        NP_{\mathrm{WCC}} = NP_{\mathrm{Turing}}.
        \]
  \item The classical complexity-theoretic separation
        $P_{\mathrm{Turing}} \neq NP_{\mathrm{Turing}}$ holds.
\end{enumerate}
Let $(D_n)_{n\ge 1}$ be any polynomial-resource WCT device family
(Definition~\ref{def:poly_resource_wct}) implementing an AGI-level
computational system on inputs of length $n$. Then there exists an
NP-complete language $L \subseteq \Sigma^*$ and a constant $k \ge 1$ such that, for infinitely many input lengths $n$, the system realized by $(D_n)$
\emph{fails} to decide $L$ on at least one input of length $n$ using only resources
obeying~\eqref{eq:poly_resource_scaling}.
Equivalently, no such WCT-bounded AGI can decide all NP-complete languages
within polynomial physical resources without violating at least one of
assumptions (i)--(iii).
\end{theorem}

\begin{proof}
Let $L_{\mathrm{SAT}}$ be any fixed NP-complete language (e.g.\ 3-SAT) over
some finite alphabet $\Sigma$. By (iii), $P_{\mathrm{Turing}} \neq
NP_{\mathrm{Turing}}$, so $L_{\mathrm{SAT}} \notin P_{\mathrm{Turing}}$. By
(ii) and the NP-equivalence, $L_{\mathrm{SAT}} \notin P_{\mathrm{WCC}}$ as
well. Unpacking the definition of $P_{\mathrm{WCC}}$, this means that for
every discrete WCC machine $W$ deciding $L_{\mathrm{SAT}}$ and for every
exponent $k \ge 1$, there exist infinitely many input lengths $n$ such that
$T_W(n) > n^{k}$.

Now let $(D_n)$ be a polynomial-resource WCT device family implementing an
AGI-level system that, by assumption (ii), corresponds to some WCC machine
$W_{\mathrm{AGI}}$ on inputs of length $n$. By
Proposition~\ref{prop:phys_step_ceiling}, there exist constants $C',d$ such
that $T_{W_{\mathrm{AGI}}}(n) \le C' n^{d}$ for all $n$, because each global
WCC update costs at most a constant number of lane updates and the total
number of lane updates is polynomially bounded.

Suppose for contradiction that this AGI decides $L_{\mathrm{SAT}}$ on all
inputs of length $n$ while respecting the polynomial resource bounds
\eqref{eq:poly_resource_scaling}. Then $W_{\mathrm{AGI}}$ decides
$L_{\mathrm{SAT}}$ with runtime $T_{W_{\mathrm{AGI}}}(n) \le C' n^{d}$, so
$L_{\mathrm{SAT}} \in P_{\mathrm{WCC}}$, and hence $L_{\mathrm{SAT}} \in
P_{\mathrm{Turing}}$ by (ii), contradicting the choice of
$L_{\mathrm{SAT}}$ and assumption (iii).

Therefore, for any polynomial-resource WCT device family $(D_n)$, there
must exist some NP-complete language $L$ (in fact any NP-complete language)
and some $k \ge 1$ such that, for infinitely many lengths $n$, the system
implemented by $(D_n)$ cannot decide $L$ on at least one input of length $n$
without exceeding the polynomial resource bounds in
\eqref{eq:poly_resource_scaling} or violating at least one of (i)--(iii).
This is precisely the claimed conditional NP-safety statement.
\end{proof}

\begin{remark}[Interpretation for AGI safety]
\label{rem:agi_safety_interpretation}
Theorem~\ref{thm:conditional_np_safety} can be read as follows. If
\emph{(a)} the WCT rails correctly bound all physically realizable
computation and \emph{(b)} the classical separation $P \neq NP$ holds, then
no AGI implemented on WCT hardware with volume, time, and energy scaling at
most polynomially in the input size can solve NP-complete problems on all
inputs within those resource bounds. In particular, such an AGI cannot
``cheat'' cryptography or other NP-hard tasks by exploiting super-Turing
parallelism or infinitely fine analog degrees of freedom: any attempt to do
so would require either super-polynomial physical resources or a violation
of the WCT rails.
\end{remark}



\section{Uniform Encodings, Decision Equivalence, and the Physical Bridge}
\label{sec:bijection}

The preceding sections establish an equivalence of \emph{language classes}
between uniform discrete local-update WCC machines and standard Turing
machines. They do not establish a pointwise bijection between every
abstract Turing configuration and every continuous wavefield.
In particular, a general polynomial-time verifier need not be in
$\NC^1$: the WCC machine simulates the full verifier computation
using polynomially many local rounds.

\subsection{Uniform input and certificate representations}
For each fixed classical deterministic machine $M$, the construction
of Section~\ref{sec:tm_wcc_equiv} supplies a polynomial-time
injective tape-to-lattice encoding $\enc_M(x)$ and a fixed-radius
local transition rule implementing one classical Turing step in a
constant number of WCC rounds. The readout of a halting state is
polynomial-time uniform. For a classical verifier $V_{\TM}(x,w)$, use
the same uniform construction on an efficiently delimited encoding
of the pair $(x,w)$.

\begin{lemma}[Preservation of verifier outcomes]
\label{lem:verifier-encoding}
For every fixed polynomial-time Turing verifier $V_{\TM}$, there
exists a uniform discrete WCC verifier $V_{\WCC}$ satisfying
\[
V_{\TM}(x,w)=1\quad\Longleftrightarrow\quad
V_{\WCC}(\enc(x,w))=1
\]
for every input and candidate witness, with polynomially bounded
encoding, time, lattice size, and decoding overhead.
\end{lemma}
\begin{proof}
Apply the two-phase radius-one Turing simulation to the fixed
verifier $V_{\TM}$. Its transition rule does not depend on $x$ or
$w$, and the input pair is encoded uniformly. The simulation
invariant preserves the simulated tape, head, and halting state,
so acceptance agrees. Polynomial runtime and tape usage of the
verifier give polynomial overhead for the WCC representation.
\end{proof}

\subsection{Decision-class equivalence}
\begin{theorem}[Uniform discrete complexity equivalence]
\label{thm:uniform-discrete-bijection}
For the uniform finite-state model of Definition~\ref{def:wcc_machine},
\[
P_{\mathrm{WCC}}^{\mathrm{disc}}=P_{\TM},
\qquad
NP_{\mathrm{WCC}}^{\mathrm{disc}}=NP_{\TM}.
\]
The equalities identify the \emph{sets of recognized languages} and
are not a claim that arbitrary encodings are bijective or that
physically constrained WCT evolution realizes every such machine.
\end{theorem}
\begin{proof}
The deterministic equality is
Corollary~\ref{cor:pwcc_equals_pturing}; the nondeterministic
verifier equality is Theorem~\ref{thm:np_equiv}.
Both rely on polynomial-time uniform input encoding and decoding.
\end{proof}

\subsection{Conditional lifting, without conflating models}
\begin{proposition}[Classical implication under an actual common-model separation]
\label{prop:conditional-classical-lift}
If a language $L$ is proved to satisfy
$L\in NP_{\mathrm{WCC}}^{\mathrm{disc}}\setminus
P_{\mathrm{WCC}}^{\mathrm{disc}}$, then $P\ne NP$.
\end{proposition}
\begin{proof}
Substitute the discrete class equalities from
Theorem~\ref{thm:uniform-discrete-bijection}.
\end{proof}

A lower bound for only \emph{physical} curvature-bounded realizations
does not satisfy the premise of this proposition. Deriving classical
$P\ne NP$ from a physical separation would require an independently
proved equality of the \emph{same} physical and discrete decision and
verifier classes, with polynomial preservation of all relevant costs.
No such physical-to-discrete resource equivalence or classical
separation is established here.


\section{Conclusion}

The principal rigorous result is the polynomial-time equivalence of
\emph{uniform discrete} WCC machines and classical Turing machines.
Fixed finite-state local rules, polynomial-size constructible lattices,
polynomial-time input encoders and output decoders, and polynomially
many global update rounds yield
\[
 P_{\mathrm{WCC}}^{\mathrm{disc}}=P,
 \qquad NP_{\mathrm{WCC}}^{\mathrm{disc}}=NP.
\]
The uniformity requirements exclude arbitrary noncomputable
preprocessing in the encoding map. Exact wave-to-symbol correspondence
requires quantization compatibility; the converse construction of an
abstract finite-state representative evolution does not establish
realization by a chosen nonlinear WCT PDE.

The physical curvature/energy channel-counting estimates are presented
conditional on their additional spectral-support and dynamical
assumptions. Neither those estimates nor the discrete Turing simulations
prove $P_{\mathrm{WCC}}^{\mathrm{phys}}=P$ or
$NP_{\mathrm{WCC}}^{\mathrm{phys}}=NP$, and no unconditional
$P\ne NP$ or universal SAT lower bound is claimed in this manuscript.

The remaining research problem is to construct and validate a precise
physical realization map with uniform polynomial resource overhead,
then separately study any model-relative physical resource lower bounds.


\section*{References}
\addcontentsline{toc}{section}{References}

\begin{enumerate}

  \item S.~A. Cook.
  ``The Complexity of Theorem-Proving Procedures.''
  \textit{Proceedings of the 3rd Annual ACM Symposium on Theory of Computing}, 1971.

  \item L.~A. Levin.
  ``Universal Sequential Search Problems.''
  \textit{Problems of Information Transmission}, 9(3):265--266, 1973.

  \item T. Baker, J. Gill, R. Solovay.
  ``Relativizations of the $P$=?$NP$ Question.''
  \textit{SIAM Journal on Computing}, 4(4):431--442, 1975.

  \item A.~A. Razborov, S. Rudich.
  ``Natural Proofs.''
  \textit{Journal of Computer and System Sciences}, 55(1):24--35, 1997.

  \item S. Aaronson, A. Wigderson.
  ``Algebrization: A New Barrier in Complexity Theory.''
  \textit{ACM Transactions on Computation Theory}, 1(1):2:1--2:54, 2009.

  \item A. Lubotzky, R. Phillips, P. Sarnak.
  ``Ramanujan Graphs.''
  \textit{Combinatorica}, 8(3):261--277, 1988.

  \item S. Arora, B. Barak.
  \textit{Computational Complexity: A Modern Approach}.
  Cambridge University Press, 2009.

  \item R.~J. Reyes.
  \textit{Structure and Derivation of Physical Constants through Wave Confinement}.
  Zenodo, 2025. \url{https://doi.org/10.5281/zenodo.15341540}.

  \item R.~J. Reyes.
  \textit{The Geometry of Resonance: Wave Confinement Theory and the Emergence of Mass, Force, and Spacetime}.
  Zenodo, 2025. \url{https://doi.org/10.5281/zenodo.15356814}.

\end{enumerate}


\end{document}
